\begin{table}[htbp]
\centering
\caption{Resumo da árvore de decisão da etapa contratual, recorte geral, avaliação in\_sample, profundidade máxima 10.}
\label{tab:tabela_compacta_treeClassification_10_profundidade_binario_recorte_geral_in_sample}
\small
\begin{tabular}{lllllll}
\toprule
\rowcolor[gray]{0.85}
Especificação & N & ROC-AUC & Renda & Desempenho & Interação & Variáveis \\
\midrule
Modelo 1 & 1.184.972 & 0,6594 & 0,5378 & 0,0431 & 0,4192 & 3 \\
\rowcolor[gray]{0.95}
Modelo 2 & 1.184.972 & 0,8322 & 0,0746 & 0,0052 & 0,0262 & 6 \\
Modelo 3 & 1.184.972 & 0,8405 & 0,0667 & 0,0041 & 0,0220 & 10 \\
\rowcolor[gray]{0.95}
Modelo 4 & 1.184.972 & 0,8453 & 0,0622 & 0,0028 & 0,0202 & 14 \\
Modelo 5 & 1.184.972 & 0,8455 & 0,0619 & 0,0025 & 0,0206 & 20 \\
\bottomrule
\end{tabular}
\par\footnotesize{Notas: A etapa contratual compara inscrições em não contratação e contratação. Renda, desempenho e interação indicam importâncias normalizadas das variáveis na árvore. A profundidade máxima parametrizada é 10.}
\par\footnotesize{Fonte: elaboração própria (2026).}
\end{table}
